\documentclass[11pt]{amsart}
\usepackage[margin=25mm]{geometry}
\usepackage{amsmath,amssymb}
\usepackage{booktabs,longtable,array}
\usepackage[hyphens]{url}
\urlstyle{tt}
\usepackage[colorlinks=true,linkcolor=blue,citecolor=blue,urlcolor=blue]{hyperref}
% The author's macros.
%
% some calligraphy
%
\newcommand\cA{{\mathcal A}}
\newcommand\cB{{\mathcal B}}
\newcommand\cC{{\mathcal C}}
\newcommand\cD{{\mathcal D}}
\newcommand\cE{{\mathcal E}}
\newcommand\cF{{\mathcal F}}
\newcommand\cG{{\mathcal G}}
\newcommand\cH{{\mathcal H}}
\newcommand\cI{{\mathcal I}}
\newcommand\cJ{{\mathcal J}}
\newcommand\cK{{\mathcal K}}
\newcommand\cL{{\mathcal L}}
\newcommand\cM{{\mathcal M}}
\newcommand\cN{{\mathcal N}}
\newcommand\cP{{\mathcal P}}
\newcommand\cQ{{\mathcal Q}}
\newcommand\cR{{\mathcal R}}
\newcommand\cS{{\mathcal S}}
\newcommand\cT{{\mathcal T}}
\newcommand\cU{{\mathcal U}}
\newcommand\cV{{\mathcal V}}

%
% fields and rings (and a semigroup)
%
\newcommand\cx{{\mathbb C}}% complexes
\newcommand\fld{{\mathbb F}}
\newcommand\flde{{\mathbb E}}
\newcommand\ints{{\mathbb Z}}
\newcommand\nn{{\mathbb N}}%non-negative integers
\newcommand\re{{\mathbb R}}%reals
\newcommand\rats{{\mathbb Q}}

\newcommand\Zu{{\mathbf u}}
\newcommand\Zv{{\mathbf v}}
\newcommand\Zw{{\mathbf w}}
\newcommand\Zy{{\mathbf y}}
\newcommand\Zz{{\mathbf z}}
\newcommand\Zx{{\mathbf x}}
\newcommand\Zj{{\mathbf j}}


%
% the really useful stuff
%
%\newcommand\opk[1]{\mathop{\hbox{\rm #1}}\nolimits} %% old version
\newcommand\opk[1]{\mathop{\mathrm{#1}}\nolimits}
\newcommand\comp[1]{{\mkern2mu\overline{\mkern-2mu#1}}}
\newcommand\diff{\mathbin{\mkern-1.5mu\setminus\mkern-1.5mu}}% for \setminus
\newcommand\gbin[2]{\genfrac{[}{]}{0pt}{}{#1}{#2}}
\newcommand\inv{^{-1}}
\newcommand\res{\mathbin{\mkern-2.0mu\restriction\mkern-2.0mu}}% ch2
\newcommand\sbs{\subseteq}
\newcommand\sps{\supseteq}
\newcommand\seq[3]{#1_{#2},\ldots,#1_{#3}}
\newcommand\supp{\opk{supp}}
\newcommand\im{\opk{im}}
\newcommand\row{\opk{row}}
\newcommand\pmat[1]{\begin{pmatrix} #1 \end{pmatrix}}
\newcommand\cprod{\mathbin{\square}}
%
% matrix theory
%
\newcommand\ip[2]{\langle#1,#2\rangle}
\newcommand\one{{\bf1}}
\newcommand\rk{\opk{rk}}
\newcommand\tr{\opk{tr}}
\newcommand\col{\opk{col}}
\newcommand\mat[3]{\mathrm{Mat}_{#1\times #2}(#3)}

%
% some group theory
%
\newcommand\aut[1]{{\rm Aut}(#1)}
\newcommand\fx[1]{{\rm fix}(#1)}% ch2
\newcommand\grp[1]{\langle #1\rangle}
\newcommand\nrml{\vartriangleleft}
\newcommand\nrmleq{\trianglelefteq}
\newcommand\alt[1]{\opk{Alt}(#1)}
\DeclareMathOperator{\Sym}{Sym}
% End of the author's macros.
% Vectors that the author's macros do not provide (as in blueprint/macros.tex).
\newcommand\Za{{\mathbf a}}
\newcommand\Zb{{\mathbf b}}
\newcommand\Ze{{\mathbf e}}
\newcommand\Zp{{\mathbf p}}

\newcommand{\cgss}{[CGSS]}
\newcommand{\file}[1]{\path{#1}}
% Lean declarations and blueprint labels, set like paths so that they break at "." and "_".
\DeclareUrlCommand\lean{\urlstyle{tt}}
\DeclareUrlCommand\lbl{\urlstyle{tt}}
% A subscript digit in a Lean name, such as the 8 of E8.
\newcommand{\lsub}[1]{\ensuremath{{}_{#1}}}
\newcommand{\errhead}[3]{\subsection*{\texorpdfstring{#1. #2}{#1}}\noindent\emph{Where.} #3\par\smallskip}
\newcommand{\printed}[1]{\noindent\emph{Printed.} #1\par\smallskip}
\newcommand{\correct}[1]{\noindent\emph{Correct.} #1\par\smallskip}
\newcommand{\missing}[1]{\noindent\emph{Missing.} #1\par\smallskip}
\newcommand{\note}[1]{\noindent\emph{Note.} #1\par\smallskip}
\newcommand{\evidence}[1]{\noindent\emph{Evidence.} #1\par\smallskip}
\newcommand{\repair}[1]{\noindent\emph{Repair.} #1\par\smallskip}
\newcommand{\status}[1]{{\raggedright\noindent\emph{Status.} #1\par}\medskip}

\title[Errors and implicit hypotheses in CGSS]{Errors and implicit hypotheses in\\
Cameron, Goethals, Seidel and Shult,\\
\emph{Line graphs, root systems, and elliptic geometry}}
\hypersetup{pdftitle={Errors and implicit hypotheses in Cameron, Goethals, Seidel and Shult}}
\author{Krystal Guo}
\date{9 October 2026}

\begin{document}
\maketitle

\section*{Scope}

This document lists the errors in the printed paper \cgss{} (J.\ Algebra \textbf{43} (1976),
305--327) and the hypotheses that its statements leave implicit, as found while formalizing the
paper. We checked every item against the published paper, and we quote the paper as printed;
page numbers are those of the journal. The permanent record
is \file{cgss/errata.md}. It also
records 19 gaps in the arguments, which are neither errors nor missing hypotheses; we do not
repeat them here, and we will cite them by their numbers there, as ``Gap 14''. Numbers in
brackets, such as [15], refer to the references of \cgss{}.

This version follows \file{cgss/errata.md}: 17 errors and 8 implicit
hypotheses. A kernel-checked certificate in Lean proves the cospectral pair behind Errors 14, 16
and 17.

Each item ends with a \emph{Status} line, which says where the correction now lives: its label in
the blueprint (the structured proofs in the style of Lamport, in \file{blueprint/}) and its Lean
declaration (in \file{LineGraphs/}), each with its file. Every numbered result of \cgss{} is
in Lean, except Remark 2.11 (Shult's characterization, which is outside the scope);
\file{dependency-diagram.md}, at the top of the repository, gives the declarations and their
files. The Lean statements state the implicit hypotheses explicitly. The statement audit
(\file{LineGraphs/AUDIT-statements.md}), which
compared the final Lean with the paper itself, found no Lean statement that says something other
than the paper without a recorded reason.

We keep the three graph conventions of the paper apart in our files: the \emph{orthogonality
graph} of Section 2 (adjacent at $90^\circ$), the \emph{represented graph} of Section 4 (inner
product 1), and the \emph{obtuse graph} before Theorem 3.8 (inner product $-1$). The quotations
keep the paper's $G$ for a graph. In our own text graphs are $X$, $Y$, $Z$; we will write $X_1$
and $X_2$ for the paper's $G_1$ and $G_2$ in Section 5, and $Y$ for the paper's graph $H$ with
$L(H)=G$ in Theorem 4.13. Strongly regular parameters are $(n,k,a,c)$. The other symbols are the
paper's, such as the parameters $v,k,b,r,\lambda$ of a design.

\section*{Summary}

In the column Status, a label is that of the blueprint, and ``Lean'' says that the corrected
statement is in Lean; each item names the declarations.

\begin{longtable}{@{}>{\raggedright\arraybackslash}p{21mm}>{\raggedright\arraybackslash}p{29mm}%
>{\raggedright\arraybackslash}p{60mm}>{\raggedright\arraybackslash}p{41mm}@{}}
\toprule
Entry & Item of \cgss & Problem & Status \\
\midrule
\endhead
Error 1 & Theorem 3.7 & $E_8$: $(C-18I)$ should be $(C-28I)$ & fixed: \lbl{thm:3.7}; Lean \\
Error 2 & Example 4.6 & coefficient $\frac13$ should be $\frac12$ & fixed: \lbl{ex:4.6}; Lean \\
Error 3 & Lemma 2.10 & entry $(3,1)$ of the block-sum matrix is 0, not 2 & fixed: \lbl{lem:2.10}; Lean \\
Error 4 & Theorem 5.8 & $s\le 2t-1$ should be $s\ge 2t-1$ & fixed: \lbl{thm:5.8}; Lean \\
Error 5 & Theorem 3.6 & coclique of $T(6)$ should be clique & fixed: before \lbl{thm:3.6}; Lean (C5) avoids it \\
Error 6 & Defs.\ 3.2, Thm.\ 2.13, proofs of 4.3, 4.4, 5.8 & labels; $\Delta$ counted by lines & fixed: \lbl{def:3.2}, \lbl{thm:2.13}; Lean \\
Error 7 & Theorem 4.13 & the count $v-2$, $v-1$ is reversed & fixed: \lbl{lem:4-LG1}; Lean \\
Error 8 & Theorem 4.4 & ``28 equiangular lines lie in $\re^7$'' is false as stated & fixed: \lbl{thm:4.4}, \lbl{lem:4-EQ1}, \lbl{lem:4-EQ2}; Lean \\
Error 9 & Section 6 & $xAJ=-nJ$ and $AJ=-dJ$ fail in general; $x\mid n$ needs $n\ge3$ & fixed: \lbl{lem:6-fourth}; Lean \\
Error 10 & Definition 3.9 (b) & $S$ should be $\Sigma$ & fixed: \lbl{def:3.9}; Lean \\
Error 11 & Remark 3.11 & four statements need qualification & fixed: \lbl{rem:3.11}, \lbl{thm:C}; Lean \\
Error 12 & Lemma 2.12 & ``the unique strongly regular graphs on 9, 15, 27 vertices'' & fixed: \lbl{lem:2.10}, \lbl{lem:2.12}; Lean \\
Error 13 & several & further slips & fixed in the blueprint \\
Error 14 & Remark 5.7 & the eigenvalues for $\alpha=vw$; the sum of two multiplicities & fixed: \lbl{rem:5.7}; Lean \\
Error 15 & Theorem 4.4 & ``$I+\frac12A$ has rank $\le7$'' fails for $(n,d)=(28,12)$ & fixed: \lbl{thm:4.4}; Lean \\
Error 16 & Theorem 5.8 & false without $b\ge v$ & fixed: \lbl{thm:5.8}; Lean \\
Error 17 & Corollary 5.9 & the exception must be unordered & fixed: \lbl{cor:5.9}; Lean \\
Impl.~hyp.~1 & Theorem 4.10; Section 1 & $X$ connected & fixed: \lbl{thm:4.10}, \lbl{lem:4-rootsystem}; Lean \\
Impl.~hyp.~2 & Theorem 4.13 & a note on how we read ``least eigenvalue $-2$''; nothing is missing & none needed: \lbl{thm:4.13}; Lean as printed \\
Impl.~hyp.~3 & Lemmas 5.4, 5.5 & no isolated vertices; an edge & fixed: \lbl{lem:5.4}, \lbl{lem:5.5}; Lean \\
Impl.~hyp.~4 & Hypothesis 2.4 & at least two lines & fixed: \lbl{hyp:2.4}, \lbl{lem:2.12}, \lbl{thm:2.13}, \lbl{thm:3.5}; Lean \\
Impl.~hyp.~5 & Lemma 5.6 & $X_2$ without isolated vertices & fixed: \lbl{lem:5.6}; Lean \\
Impl.~hyp.~6 & Section 5 & meaning of ``characterized by its spectrum'' & fixed: \lbl{def:flag}; Lean \\
Impl.~hyp.~7 & Definition 2.2 & maximal within the span of the lines & fixed: \lbl{def:2.2}; Lean \\
Impl.~hyp.~8 & Theorem 3.12 & either part of the partition may be empty & fixed: \lbl{thm:3.12}; Lean \\
\bottomrule
\end{longtable}

\section{Errors in the printed paper}

\errhead{Error 1}{Theorem 3.7, the constant for $E_8$}{Theorem 3.7; p.~312.}
\printed{``$E_8$: $(C+2I)(C-18I)=0$''.}
\correct{$(C+2I)(C-28I)=0$. The bound 28 on the valency in Theorem 4.3 is this constant (Gap 14);
with the printed 18, the same argument would give 18, which the example in Theorem 4.3
contradicts.}
\evidence{Take one vector of squared norm 2 on each of the 120 lines of $E_8$. These vectors form
a tight frame in $\re^8$, $\sum\Zv\Zv^T=30I$, so the Gram matrix $2I+C$ satisfies
$(2I+C)^2=30(2I+C)$, and $C$ has eigenvalues 28 and $-2$, with multiplicities 8 and 112.
Equivalently, each line of $E_8$ is at $60^\circ$ to $56=2\cdot28$ others, and the trace argument
that the paper gives for $D_n$ yields $\gamma_1+\dots+\gamma_8=224$, so each $\gamma_i$ is 28. We
checked the eigenvalues numerically for $E_6$, $E_7$ and $E_8$; the printed values 10 and 16 for
$E_6$ and $E_7$ are correct.}
\repair{We derive all five constants from the frame identity. Let $N$ be the number of lines and
$\ell$ the rank. Over one vector per line, $\sum\Zv\Zv^T=\frac{2N}{\ell}I$ on the span of the roots
(computation C6); in coordinates, this is $\frac{2N}{\ell}$ times the orthogonal projection onto
the span, since the models of $A_n$, $E_7$, $E_6$ and the first model of $E_8$ do not span their
coordinate spaces. The second constant is then $\frac{2N}{\ell}-2$, which is 28 for $E_8$, where
$\frac N\ell=15$. The diagonal of $C^2$ is LN6 of \file{cgss/prerequisites.md} (computation C20).
The same derivation is in \file{cgss/proofs/remark311/remark311.tex}.}
\status{Fixed. Blueprint: Theorem 3.7, \lbl{thm:3.7}, through the frame identity within the span,
\lbl{lem:sec3-frame}, and C6, \lbl{comp:C6} (\file{blueprint/sec3-rootsystems.tex}). Lean:
\lean{LineSystem.signedAdj_quadratic_E}\lsub8, and the same name ending in \lean{_An}, \lean{_Dn},
\lean{_E}\lsub7 and \lean{_E}\lsub6 for the other types, all from \lean{LineSystem.frame_identity}
(\file{LineGraphs/Computations/C6.lean}); C6 is \lean{Computations.C6.coordinate_frame_matrices}.}

\errhead{Error 2}{Example 4.6, the Schl\"afli graph}{Example 4.6; p.~319.}
\printed{``$\{\frac13\Sigma e_k-e_i-e_7;\ i\in T(6,1)\}$''.}
\correct{$\frac12\sum_{k=1}^8\Ze_k-\Ze_i-\Ze_7$, as in the third family.}
\evidence{With $\frac13$ these vectors have squared norm $14/9$, so they do not lie in the model
(ii) of $E_8$. With $\frac12$ the 27 vectors lie in $E_8$ and represent a strongly regular graph
with parameters $(27,16,10,8)$ (computation C9).}
\repair{We use $\frac12$; computation C9 checks the vectors and the parameters.}
\status{Fixed. Blueprint: Example 4.6, \lbl{ex:4.6}, with C9, \lbl{comp:C9}
(\file{blueprint/sec4-leasteigenvalue.tex}). Lean:
\lean{SimpleGraph.schlafli_isRepresentedIn_E}\lsub8 and \lean{SimpleGraph.schlafli_isSRGWith}
(\file{LineGraphs/Computations/C9Schlafli.lean}).}

\errhead{Error 3}{Lemma 2.10, the matrix of block sums}{Proof of Lemma 2.10; p.~309.}
\printed{The third row of the $5\times5$ matrix of block sums is
``$2\quad0\quad k^2+k\quad k^2-k\quad k^2-k$''.}
\correct{Its first entry is 0. The block of the first matrix that it sums, the inner products of
the vertices of $K$ with $u$, is 0, since the vertices of $K$ are adjacent to $u$; as printed, the
second matrix is not even symmetric.}
\evidence{Symmetry, and a direct computation of the block sums of the first matrix. The
determinant that the proof uses is that of the matrix with its first row and column deleted, which
does not contain the wrong entry; it is $2k^5(5-k)$, as stated. Thus the argument is unaffected.}
\repair{The blueprint recomputes the block sums, with 0 in position $(3,1)$; computation C1 is the
determinant.}
\status{Fixed. Blueprint: Lemma 2.10, \lbl{lem:2.10}, step \lbl{210gram}, which computes the
block sums in its substeps \lbl{210blocks} and \lbl{210sums} and corrects the entry in its closing
step, and C1, step \lbl{210det} (\lbl{comp:C1}), in \file{blueprint/sec2-lines.tex}.
Lean: \lean{Computations.C1.blockSumMatrix}, with the entry 0, is the Gram matrix of the five
block sums by \lean{LineSystem.Hyp24.gram_blockSumVectors}
(\file{LineGraphs/LineSystem/OrthGraphGram.lean}); C1 is \lean{Computations.C1.det_blockSumMinor}
(\file{LineGraphs/Computations/C1.lean}).}

\errhead{Error 4}{Theorem 5.8, the direction of an inequality}{Proof of Theorem 5.8; p.~323.}
\printed{``$k\leqslant v$ implies $s\leqslant 2t-1$''.}
\correct{$s\ge 2t-1$, as in the statement, case (iii).}
\evidence{Here $k=st$ and $v=\frac12 s(s+1)$ with $s,t>0$ (the remark after Theorem 5.8), so
$k\le v$ if and only if $t\le\frac12(s+1)$, that is, $s\ge 2t-1$.}
\repair{The blueprint of Theorem 5.8 proves $s\ge2t-1$ from $k\le v$; computation C11.}
\status{Fixed. Blueprint: Theorem 5.8, \lbl{thm:5.8}, step \lbl{s558nbineq}
(\file{blueprint/sec5-spectral.tex}). Lean: \lean{Computations.C11.two_mul_sub_one_le}
(\file{LineGraphs/Computations/C11.lean}), used by Theorem 5.8, \lean{Design.flagGraph_cospectral}
(\file{LineGraphs/LineGraphSpectra/Theorem58.lean}).}

\errhead{Error 5}{Theorem 3.6, clique and coclique}{Proof of Theorem 3.6; p.~312.}
\printed{``$A_7\subset E_7$ since a 5-coclique is imbeddable in the triangular graph on 6
symbols''.}
\correct{The argument mixes the two graphs on $\Gamma_a$. In the graph of Section 2, where
adjacency is orthogonality, $\Gamma_a$ of $A_7$ is a coclique on 5 vertices, and $\Gamma_a$ of
$E_7$ is the complement of $T(6)$, with parameters $(15,6,1,3)$. The inclusion needs a 5-coclique
of the complement of $T(6)$, that is, a 5-clique of $T(6)$: the five pairs through a fixed symbol.
$T(6)$ itself has no coclique on more than 3 vertices. The conclusion stands.}
\evidence{The definitions; computation C5.}
\repair{Computation C5 proves $A_7\subset E_7$ by an explicit inclusion of coordinates, and does
not need the 5-clique.}
\status{Fixed. Blueprint: a remark just before Theorem 3.6 (\lbl{thm:3.6}) records the correction,
and C5, \lbl{comp:C5}, proves the inclusion (\file{blueprint/sec3-rootsystems.tex}). Lean:
\lean{Computations.C5.inclusions_and_obstructions} (\file{LineGraphs/Computations/C5.lean}), which
Theorem 3.6, \lean{LineSystem.classification_of_isMaximal}, uses.}

\errhead{Error 6}{Minor inconsistencies of labels and counts}{Definition 3.2 (p.~311), Theorem 2.13
(p.~310), the proofs of Theorems 4.3 and 4.4 (pp.~317--318) and of Theorem 5.8 (p.~323).}
\printed{Definition 3.2 labels its second model ``(i)''; later references call it (ii): ``We use
Definition 3.2(ii) for $E_8$'' (p.~318). Theorem 2.13 gives ``$|\Delta|=36$'', ``$|\Delta|=15$'',
``$|\Delta|=6$''. The proofs of Theorems 4.3 and 4.4 cite ``Lemma 3.7'' for Theorem 3.7, and the
proof of Theorem 5.8 cites ``Theorem 5.6'' for Lemma 5.6.}
\correct{The second model is (ii). The numbers of Theorem 2.13 count lines: as defined in
Hypothesis 2.4, $\Delta$ is a set of vectors closed under negation, with 72, 30 and 12 elements.
The line counts are the ones for which $3+3|\Gamma_a|+|\Delta|$ is 120, 63 and 36. The citations
are of Theorem 3.7 and Lemma 5.6.}
\evidence{The printed text.}
\repair{The blueprint uses the correct labels and names the three models of $E_8$ (i), (ii),
(iii). In Theorem 2.13 it writes $|\Delta|/2$ for the number of lines; in Lean, $\Delta$ is a set
of vectors.}
\status{Fixed. Blueprint: Definition 3.2, \lbl{def:3.2} (\file{blueprint/sec3-rootsystems.tex});
Theorem 2.13, \lbl{thm:2.13}, step \lbl{213exceptional} (\file{blueprint/sec2-lines.tex}); the
citation of Theorem 3.7 for the valencies, in Lemma \lbl{lem:4-valency}
(\file{blueprint/sec4-leasteigenvalue.tex}); and that of Lemma 5.6 in Theorem 5.8, \lbl{thm:5.8}
(\file{blueprint/sec5-spectral.tex}). Lean: the three models are
\lean{LineSystem.E}\lsub8\lean{viaA}\lsub8, \lean{LineSystem.E}\lsub8\lean{viaD}\lsub8 and
\lean{LineSystem.E}\lsub8\lean{viaSteiner} (\file{LineGraphs/LineSystem/Models.lean}), and
\lean{LineSystem.five_possibilities'} gives $|\Delta|/2=36,15,6$
(\file{LineGraphs/LineSystem/ClassificationModels.lean}).}

\errhead{Error 7}{Theorem 4.13, a reversed count}{Proof of Theorem 4.13; p.~320.}
\printed{``If $\{i,j\}$ and $\{i,k\}$ are edges of $H$, then the number of vertices of $L(H)$
adjacent to both is $v-2$ or $v-1$ according as $\{j,k\}$ is an edge or not.''}
\correct{It is $v-1$ if $\{j,k\}$ is an edge and $v-2$ otherwise.}
\evidence{Here $Y$, the paper's $H$, is regular of valency $v$. The common neighbours of $\{i,j\}$
and $\{i,k\}$ in $L(Y)$ are the $v-2$ other edges on $i$, together with $\{j,k\}$ when it is an
edge. The argument uses only that the two values differ, so it is unaffected.}
\repair{The line graph case LG1 of \file{cgss/prerequisites.md}, \S7.2, which the blueprint
proves for Theorem 4.13, uses the corrected count.}
\status{Fixed. Blueprint: Lemma (the line graph case, LG1), \lbl{lem:4-LG1}, used by Theorem 4.13,
\lbl{thm:4.13} (\file{blueprint/sec4-leasteigenvalue.tex}). Lean:
\lean{SimpleGraph.lineGraph_isSRGWith_cases} (\file{LineGraphs/Seidel/LineGraphCases.lean}), which
the proof of \lean{SimpleGraph.seidel} uses.}

\errhead{Error 8}{Theorem 4.4, 28 equiangular lines}{Proof of Theorem 4.4; p.~318.}
\printed{``But any set of 28 equiangular lines lies in $\re^7$, hence the multiplicity of 0 as an
eigenvalue actually is $\geqslant n-7$. This implies that $I+\frac12A$ has rank $\leqslant 7$,
hence the set of vectors we started with is contained in $E_7$ or $D_7$.''}
\correct{The first sentence is false as stated: 28 lines at angle $\arccos\frac13$ can span
$\re^{15}$. It is true for lines in $\re^m$ with $m\le14$, in particular in $\re^8$, which is the
case used, but [15] (Lemmens and Seidel, \emph{Equiangular lines}) does not state it (Gap 7). The
last clause should also allow $A_7$ and the smaller systems; since $A_7\subset E_7$, every
indecomposable star-closed set of rank at most 7 lies in $E_7$ or $D_7$, and nothing changes.}
\evidence{A single full pillar in the sense of [15] gives 28 lines at $\arccos\frac13$ spanning
$\re^{15}$ (\file{cgss/prerequisites.md}, \S7.1).}
\repair{EQ1 to EQ3 of \file{cgss/prerequisites.md}: the relative bound of [15], and the cone
construction through Theorem 3.5 (decision 1 of \file{cgss/prerequisites.md}, \S13).}
\status{Fixed. Blueprint: Theorem 4.4, \lbl{thm:4.4}, through the relative bound EQ1,
\lbl{lem:4-EQ1}, and the lines spanning eight dimensions, EQ2, \lbl{lem:4-EQ2}
(\file{blueprint/sec4-leasteigenvalue.tex}). Lean: \lean{InnerProductSpace.card_le_relativeBound}
(\file{LineGraphs/Representation/Theorem44.lean}) and
\lean{InnerProductSpace.card_le_fourteen_of_finrank_eq_eight}
(\file{LineGraphs/Spectral/RelativeBoundFourteen.lean}); Theorem 4.4 is
\lean{SimpleGraph.IsRepresentedIn.card_le_28_of_regular} and
\lean{SimpleGraph.IsRepresentedIn.degree_le_16_of_regular}
(\file{LineGraphs/Representation/Theorem44Unconditional.lean}).}

\errhead{Error 9}{Section 6, the row sums of $A$}{Section 6, before Theorem 6.2; p.~326.}
\printed{``$(xA+nI)(xA-\frac12n(n-2)I)=0$, $xAJ=-nJ$, hence $x$ divides $n$'' and
``$(A+dI)(A-\frac12d(n-2)I)=0$, $AJ=-dJ$''.}
\correct{The quadratic equations hold; the equations $xAJ=-nJ$ and $AJ=-dJ$ do not hold in
general. Since $PP^T=nI+xA$, the equation $xAJ=-nJ$ says that $PP^T\Zj=0$, that is, $P^T\Zj=0$.
The entry of $P^T\Zj$ at a column $c$ is $\frac12(\sigma(c)^2-n)$, where $\sigma(c)$ is the sum of
column $c$ of $H_n$. Thus $xAJ=-nJ$ holds if and only if every column sum of $H_n$ is $\pm\sqrt n$.
This depends on the signs of the rows of $H_n$, and it fails for every choice of signs unless $n$
is a square. The conclusion that $x$ divides $n$ is true for another reason: as the paper says
next, $d^2=(n/x)^2$ is the number of fourth rows that agree with three given rows in
$\frac18(n\pm x)$ places, so $n/x$ is a rational number whose square is an integer. (For a fixed
row $i$, the $n-1$ vectors $H_i\circ H_j$ form an orthogonal basis of $\Zj^\perp$, and expanding
$H_k\circ H_l$ in this basis gives the count.) The claim needs $n\ge3$. For $n=1$ and $n=2$ the
matrix has fewer than four rows, so Hypothesis 6.1 holds for every $x$, and $x$ need not divide
$n$ (take $n=2$ and $x=3$); there is no Hadamard matrix of order 3. The assumption $x>0$ costs
nothing: for $n\ge4$ and $x=0$ we would have $PP^T=nI$, of rank $\binom n2>n$, whereas $P$ has
only $n$ columns.}
\evidence{For the Paley matrix $H_{12}$, the row sums of $xA$ are $-60$ and $36$, and no choice of
row signs makes every column sum $\pm\sqrt{12}$. For the Sylvester matrix of order 8, every row
sum of $xA$ is $24=\frac12n(n-2)$, so $\Zj$ is an eigenvector of $PP^T$ for $\frac12n^2$, not
for 0 (NumPy, scratch). The same check confirms Hypothesis 6.1 for $H_{12}$ with $x=4$. Theorem 6.2
is unaffected, since its proof uses only the quadratic equation, the fact that the lines lie in
$\Zj^\perp$, and the orthogonality of the lines through a fixed row.}
\repair{The blueprint states the quadratic equation, and derives $x\mid n$ from the count of
fourth rows; computation C16. A shorter route: for $n>2$ the eigenvalue $-n/x$ of the integer
matrix $C_H$ (the paper's $A$) has multiplicity $\frac12(n-1)(n-2)>0$, so it is an integer by the
rational root theorem (\file{cgss/proofs/spectraltools/spectraltools.tex}, Theorem 9.6, used in its
Theorem 10.4).}
\status{Fixed. Blueprint: the background of Section 6, the quadratic equation,
\lbl{lem:6-quadratic}, and the fourth rows, \lbl{lem:6-fourth}, with the remark after it
(\file{blueprint/sec6-hadamard.tex}). Lean: \lean{Matrix.IsHadamard.hadamardSignedAdj_quadratic}
(\file{LineGraphs/Hadamard/Agreement.lean}) and \lean{Matrix.IsHadamard.FourRowAgreement.dvd_card}
(\file{LineGraphs/Hadamard/FourthRows.lean}), which assumes $n\ge3$ and $x>0$. The two false
identities have no Lean statement; the statement audit records them as rightly absent.}

\errhead{Error 10}{Definition 3.9 (b), $S$ for $\Sigma$}{Definition 3.9 (b); p.~314.}
\printed{``$\forall_{r\in\Sigma}\forall_{\lambda\in\re}((\lambda r\in S)\Rightarrow(\lambda\in\{1,-1\}))$''.}
\correct{$\lambda r\in\Sigma$. The letter $S$ names the set of lines of Section 2.}
\evidence{The printed text, and the sense of the definition. Read with $\Sigma$, condition (b) also
excludes 0 from $\Sigma$ (if $0\in\Sigma$, then $\lambda0\in\Sigma$ for $\lambda=2$).}
\repair{Section 2 of \file{cgss/proofs/remark311/remark311.tex}, and item 2 of its section
``Interpretations and corrections of the sources''.}
\status{Fixed. Blueprint: Definition 3.9, \lbl{def:3.9} (\file{blueprint/sec3-remark311.tex}).
Lean: \lean{IsRootSet} (\file{LineGraphs/RootSystems/Defs.lean}), whose field \lean{reduced} is
condition (b), read with $\Sigma$.}

\errhead{Error 11}{Remark 3.11, statements that need qualification}{Remark 3.11; p.~314.}
\printed{``Let $r,s\in\Sigma$ be not perpendicular. Since $4(r,s)^2/(r,r)(s,s)$ is a positive
integer $<4$, \dots''; ``the subset of a Root System consisting of vectors of any fixed length is
itself a Root System, necessarily one of those described by Theorem 3.10''; ``$B_n$ and $C_n$: a
set of $n^2$ lines in $\re^n$ with angles $90^\circ$, $60^\circ$, $45^\circ$, and equation
$(C+2I)(C-2(n-1)I)=0$'', and the similar lines for $F_4$ and $G_2$.}
\correct{(1) The roots $r$ and $s$ must also be non-proportional; for $s=\pm r$ the quotient is 4.
(2) The roots of one length form a Root System in their own span, which need not be the whole space
and need not be indecomposable: the long roots of $B_2$ form $A_1\oplus A_1$, and the roots of
length 1 in $\{\pm\Zu_1,\pm2\Zu_2\}$ span a line. (3) The matrix $C$ is defined before Theorem 3.7
only for vectors of one length. For two root lengths we take unit vectors on the lines, as the
columns of a matrix $U$, and put $C=2U^TU-2I$; then the three printed equations hold. (4) The
listed angles are the angles allowed: for $B_2=C_2$ the angle $60^\circ$ does not occur.}
\evidence{\file{cgss/proofs/remark311/remark311.tex}, \S~``Interpretations and corrections of the
sources'', items 3 to 6, and its Lemma ``angles'' and Theorem C, with their checks.}
\repair{\file{cgss/proofs/remark311/remark311.tex} (decision 5 of \file{cgss/prerequisites.md},
\S13); the blueprint of Remark 3.11 follows it.}
\status{Fixed. Blueprint: Remark 3.11, \lbl{rem:3.11}, through Lemma (angles and fixed-length
subsystems), \lbl{lem:angles}, and Theorem C, \lbl{thm:C} (\file{blueprint/sec3-remark311.tex}).
Lean: \lean{IsRootSet.remark_lines} (\file{LineGraphs/RootSetClassification/RemarkLines.lean}),
with \lean{RootSetModel.Bn_Cn_same_lines} and the equations
\lean{RootSetModel.signedAdj_quadratic_Bn}, \lean{_Cn}, \lean{_F}\lsub4 and \lean{_G}\lsub2; the
subsets of one length are \lean{IsRootSet.fixedLength_isRootSet}.}

\errhead{Error 12}{Lemma 2.12, ``the unique strongly regular graphs on 9, 15, 27 vertices''}{Lemma
2.12; p.~310.}
\printed{``\dots the void graph on $n$ vertices, the graph consisting of $n$ triangles meeting at
one vertex, the unique strongly regular graphs on 9, 15, 27 vertices.''}
\correct{There are other strongly regular graphs on these numbers of vertices: $T(6)$ on 15 and
the Schl\"afli graph on 27 (the complements of the intended graphs), and the degenerate $3K_3$ and
$K_{3,3,3}$ on 9. The intended graphs are those of Lemma 2.10, which are strongly regular with
parameters $(9,4,1,2)$, $(15,6,1,3)$ and $(27,10,1,5)$. The paper proves only that they are
regular. Strong regularity takes two further observations, in the notation of Lemma 2.10: adjacent
vertices have exactly one common neighbour, the third vertex of their triangle; and nonadjacent
vertices $u$, $v$ have the $k$ common neighbours in $K$.}
\evidence{$T(6)$ has parameters $(15,8,4,4)$ and the Schl\"afli graph $(27,16,10,8)$.}
\repair{The blueprint names the graphs by their parameters, or as the collinearity graphs of the
generalized quadrangles of orders $(2,1)$, $(2,2)$, $(2,4)$ (GQ1 and GQ2 of
\file{cgss/prerequisites.md}).}
\status{Fixed. Blueprint: Lemma 2.10, \lbl{lem:2.10}, proves strong regularity with parameters
$(6k-3,2k,1,k)$, and Lemma 2.12, \lbl{lem:2.12}, lists the graphs by their parameters
(\file{blueprint/sec2-lines.tex}). Lean: \lean{LineSystem.Hyp24.orthGraph_isSRGWith}
(\file{LineGraphs/LineSystem/OrthGraphParameters.lean}) and \lean{LineSystem.Hyp24.orthGraph_cases}
(\file{LineGraphs/LineSystem/OrthGraphCases.lean}).}

\errhead{Error 13}{Further slips}{Various pages.}
\printed{(1) Lemma 3.4: ``$|A|=\frac12n(n+1)$'' for $|A_n|$ (p.~311). (2) Definition 3.2 (ii):
``$\epsilon_i=\pm1$'' appears twice in the display (p.~311). (3) Theorem 3.12: ``two lines in
different subsets are at angles $90^\circ$ and $45^\circ$'' for ``$90^\circ$ or $45^\circ$''
(p.~315). (4) Section 4: ``is and orthonormal basis'' (p.~315); ``by delecting a one-factor''
(p.~316); ``whose lines from a subset of $D_n$'' (p.~316); and Corollary 5.10: ``The graphs ar
switching equivalent'' (p.~324). (5) The representation of the generalized line graphs, and the
proof of Theorem 4.2: ``$1\leqslant\alpha\leqslant\alpha_i$'' for $1\le\alpha\le a_i$ (pp.~316
and 317). (6) Proof of Theorem 4.4: ``The projections onto the lines spanned by these vectors are
elements of $\re^{28}$ and have Gram matrix $(1-\frac{d+2}{2n})^2I+\cdots$'' (p.~317). (7) Proof
of Theorem 4.4: ``for each eigenvalue $\alpha\neq d$ of $A$'' (p.~318). (8) Remark 5.7: ``two
eigenvalues $\delta:=\gamma-\frac12(v+w)$ and $\delta:=\gamma'-\frac12(v+w)$'' (p.~322).}
\correct{Items (1) to (5) are misprints with evident corrections. In (6), the vectors have squared
norm $1-\frac{d+2}{2n}$, and the displayed matrix is the Gram matrix of the rank-one matrices
$\Zp\Zp^T$, not of the projections, whose Gram matrix is the displayed one divided by
$(1-\frac{d+2}{2n})^2$; the conclusion is the same. In (7), read ``for each eigenvalue $\alpha$ of
$A$ with an eigenvector orthogonal to $\Zj$''. The two readings differ only when the graph is
disconnected; then $d$ is also an eigenvalue on $\Zj^\perp$, the inequality holds for it, and the
sum over $\alpha_1,\dots,\alpha_7$ needs it. In (8), the second is $\delta':=\gamma'-\frac12(v+w)$.}
\evidence{The printed text. None of these affects an argument.}
\repair{The blueprint uses the corrected forms.}
\status{Fixed in the blueprint: (1) Lemma 3.4, \lbl{lem:3.4}; (2) Definition 3.2, \lbl{def:3.2};
(3) Theorem 3.12, \lbl{thm:3.12}, whose Lean statement,
\lean{LineSystem.subset_rootSet_iff_bipartition}, also reads ``or''; (4) Theorem 4.2,
\lbl{thm:4.2}, for ``form''; (5) Lemma (the representation of a generalized line graph),
\lbl{lem:4-glgrep}; (6) and (7) Theorem 4.4, \lbl{thm:4.4}; (8) Remark 5.7, \lbl{rem:5.7}. The
Lean statements are those of the corrected forms; for instance, \lean{LineSystem.numLines_An}
counts the lines of $A_n$.}

\errhead{Error 14}{Remark 5.7, the eigenvalues of $A(X_2)$}{Remark 5.7; pp.~322--323.}
\printed{``To each eigenvalue $\alpha\neq0$ of $B_1B_1^T$ there correspond \dots\ two eigenvalues
$\delta:=\gamma-\frac12(v+w)$ and $\delta:=\gamma'-\frac12(v+w)$ of $A(G_2)$'' (the second
$\delta$ is the misprint of Error 13, item 8) and ``For an eigenvalue 0 of $A(G_1)$ the
corresponding eigenvalues of $A(G_2)$ are $\pm\frac12(v-w)$. Their multiplicities need not be
equal, add up to $q-p$''.}
\correct{In the paper's notation, the parts of $X_1$ have sizes $p\le q$ and valencies $v$ and $w$,
and $\rho=\rk B_1$. (a) For $\alpha=vw$ the roots of $(\gamma-v)(\gamma-w)=\alpha$ are $v+w$ and
0, and 0 is not an eigenvalue of $M_2M_2^T$, which is nonsingular; thus only $\frac12(v+w)$
occurs, not $-\frac12(v+w)$. (b) The multiplicities of $\pm\frac12(v-w)$ are $p-\rho$ and $q-\rho$; their
\emph{difference} is $q-p$, and their sum is $p+q-2\rho$.}
\evidence{Take $X_1=K_{6,10}$ and $X_2=T(6)$. Then $L(K_{6,10})$ and $L(T(6))$ both have spectrum
$14^1\,8^5\,4^9\,(-2)^{45}$, $X_2$ is 8-regular and not bipartite, and $|X_1|=|X_2|+1$; the
eigenvalues $\pm2$ of $T(6)$ have multiplicities 5 and 9, which add up to $14\ne q-p=4$. This is
the situation of Corollary 5.9 and Remark 5.12, where $B_1=J$ has rank 1. We checked it by
hand, and a kernel-checked certificate also proves it (see the Status). Theorem 5.8 does not use
these statements.}
\repair{The spectral toolkit, Corollary 6.8 of \file{cgss/proofs/spectraltools/spectraltools.tex}, states the
correct spectrum, and the blueprint of Remark 5.7 uses it.}
\status{Fixed. Blueprint: Remark 5.7, \lbl{rem:5.7} (\file{blueprint/sec5-spectral.tex}), which
gives the spectrum of $A(X_2)$, with the multiplicities $p-\rho$ and $q-\rho$. Lean:
\lean{SimpleGraph.lineGraph_cospectral_bipartite_spectrum_of_not_isBipartite}
(\file{LineGraphs/LineGraphSpectra/Nonbipartite.lean}). The evidence is kernel-checked: the
spectra of $L(K_{10,6})$, the flag graph of the trivial design $(6,6,10,10,10)$, of
$L(L(K_6))=L(T(6))$ and of $T(6)$ are \lean{Computations.TrivialDesignException.spectrum_flagGraph},
\lean{.spectrum_double_lineGraph} and \lean{.spectrum_triangular}, and the cospectrality is
\lean{.flagGraph_cospectral_double_lineGraph}
(\file{LineGraphs/Computations/TrivialDesignException.lean}).}

\errhead{Error 15}{Theorem 4.4, the rank of $I+\frac12A$ when $d=12$}{Proof of Theorem 4.4; p.~318.}
\printed{The matrix $I+\frac23A-\frac13(J-I)$ ``has the simple eigenvalue
$1+\frac23d-\frac13(n-1)>0$'', and later ``This implies that $I+\frac12A$ has rank $\leqslant7$,
hence the set of vectors we started with is contained in $E_7$ or $D_7$.''}
\correct{For $n=28$ the eigenvalue on $\Zj$ is $\frac13(2d-24)$, which is 0 when $d=12$, and
$(n,d)=(28,12)$ is the case of $T(8)$ and the Chang graphs, which lie in $E_8$. For $T(8)$, with
spectrum $12^1\,4^7\,(-2)^{20}$, the matrix $I+\frac23A-\frac13(J-I)$ has eigenvalues $0$, $4^7$,
$0^{20}$, and so rank 7, while $I+\frac12A$ has eigenvalues $7$, $3^7$, $0^{20}$, and so rank 8.
Thus the sentence is false for $d=12$. The conclusion $d\le16$ stands, since $12\le16$.}
\evidence{We checked it by hand from the spectrum of $T(8)$.}
\repair{The blueprint of Theorem 4.4 splits into $d\le12$, where there is nothing to prove for the
bound, and $d>12$, where the eigenvalue on $\Zj$ is positive and the rank argument goes through.}
\status{Fixed. Blueprint: Theorem 4.4, \lbl{thm:4.4}, step \lbl{s44rankE}, with the case $d\le12$
in step \lbl{s44big} (\file{blueprint/sec4-leasteigenvalue.tex}). Lean:
\lean{SimpleGraph.IsRepresentedIn.degree_le_16_of_regular}
(\file{LineGraphs/Representation/Theorem44Unconditional.lean}); that the case $(28,12)$ occurs is
part of \lean{SimpleGraph.regular_28_12_and_27_16_isRepresentedIn_E}\lsub8
(\file{LineGraphs/Representation/Theorem44.lean}).}

\errhead{Error 16}{Theorem 5.8 needs $b\ge v$}{Theorem 5.8 and the remark after it; pp.~323--324.}
\printed{``Let $D_1$ be a 2-design with parameters $v,k,b,r,\lambda$ \dots\ Then one of the
following holds'', with the cases (i) ``$v+b\leqslant9$, \dots'', (ii) ``$G_2=F(D_2)$, \dots''
and (iii) ``$v=\frac12s(s+1)$, $k=st$, $b=\frac12(s+1)(s+2)$, \dots''.}
\correct{The statement fails without $b\ge v$. Take the trivial design $v=k=10$,
$b=r=\lambda=6$: its flag graph $L(K_{10,6})$ is cospectral with $L(T(6))$ (Error 14). Case (i)
fails, since $v+b=16$; (ii) fails, since the only design with these parameters has flag graph
$L(K_{10,6})$, and $T(6)$ is not bipartite; and (iii) fails, since $v=10$ forces $s=4$, and then
$b=15$. The dual design ($v=6$, $b=10$) satisfies (iii) with $s=3$, $t=2$. The paper's remark
after the theorem ($s,t>0$ when $\lambda<r$; otherwise pass to the dual) is the intended
hypothesis, and every corollary meets it.}
\evidence{We checked it by hand from the spectra in Error 14, and a kernel-checked certificate
proves that $L(K_{10,6})$ and $L(T(6))=L(L(K_6))$ are cospectral and not isomorphic.}
\repair{The blueprint of Theorem 5.8 assumes $b\ge v$.}
\status{Fixed. Blueprint: Theorem 5.8, \lbl{thm:5.8} (\file{blueprint/sec5-spectral.tex}). Lean:
\lean{Design.flagGraph_cospectral}, with the hypothesis \lean{hvb} that $v\le b$
(\file{LineGraphs/LineGraphSpectra/Theorem58.lean}). The evidence:
\lean{Computations.TrivialDesignException.lineGraph_completeBipartite_cospectral_double_lineGraph}
and \lean{.flagGraph_not_nonempty_iso_double_lineGraph}, the latter for the flag graph of the dual
design $(6,6,10,10,10)$, which is $L(K_{10,6})$ again
(\file{LineGraphs/Computations/TrivialDesignException.lean}).}

\errhead{Error 17}{Corollary 5.9, the exception is unordered}{Corollary 5.9; p.~324.}
\printed{``the graph $L(K_{m,n})$ is characterized by its spectrum, unless $m=2t^2+t$,
$n=2t^2-t$, and there exists a symmetric Hadamard matrix with constant diagonal of order
$4t^2$.''}
\correct{The exception must read $\{m,n\}=\{2t^2+t,2t^2-t\}$: for $(m,n)=(6,10)$ the printed
condition does not apply, but $L(K_{6,10})=L(K_{10,6})$ is not characterized (it is cospectral
with $L(T(6))$; $t=2$).}
\evidence{The pair of Error 14. A kernel-checked certificate proves that the
flag graph $L(K_{10,6})$ of the trivial design $(6,6,10,10,10)$ is cospectral with
$L(L(K_6))=L(T(6))$, is not isomorphic to it, and is not characterized by its spectrum in the
sense of Implicit hypothesis 6.}
\repair{The blueprint states Corollary 5.9 for the unordered pair $\{m,n\}$, and reduces it to
$m\ge n$ by symmetry.}
\status{Fixed. Blueprint: Corollary 5.9, \lbl{cor:5.9}, step \lbl{s559ordered}
(\file{blueprint/sec5-spectral.tex}). Lean: \lean{SimpleGraph.lineGraph_completeBipartite_characterized},
whose exception \lean{SimpleGraph.HasCompleteBipartiteHadamardException} is symmetric in $m$ and
$n$, through \lean{SimpleGraph.lineGraph_completeBipartite_characterized_of_le}
(\file{LineGraphs/LineGraphSpectra/Corollaries.lean}). The evidence:
\lean{Computations.TrivialDesignException.flagGraph_cospectral_double_lineGraph},
\lean{.flagGraph_not_nonempty_iso_double_lineGraph} and \lean{.not_flagGraphCharacterized}
(\file{LineGraphs/Computations/TrivialDesignException.lean}).}

\section{Implicit hypotheses}

The Lean statements state each of these explicitly. Implicit hypothesis 2 is a note on how we
read a phrase of the paper, not a hypothesis that the paper leaves out; it keeps its number, which
other files cite.

\errhead{Implicit hypothesis 1}{Theorem 4.10 needs $X$ to be connected}{Theorem 4.10; p.~319. The
claim of the introduction that the lines of such a graph lie in one of the root systems needs the
condition too; p.~306.}
\printed{``If $G$ is a graph with least eigenvalue $-2$, then either (i) $G$ is a generalized line
graph; or (ii) $G$ is represented by a subset of $E_8$.'' In the introduction: ``Any graph with
least eigenvalue $-2$ may be viewed as a set of lines at $60^\circ$ and $90^\circ$. As a
consequence of the main theorem, this set is contained in one of the root systems.''}
\missing{$X$ is connected.}
\evidence{Two disjoint copies of the Petersen graph have least eigenvalue $-2$ and do not form a
generalized line graph. Their matrix $2I+A$ has rank $12>8$, so they are not represented in $E_8$
either. Thus their lines lie in none of the root systems $A_n$, $D_n$, $E_6$, $E_7$, $E_8$: subsets
of $A_n$ and $D_n$ represent generalized line graphs (Theorems 4.1 and 4.2), and $E_6$, $E_7$ and
$E_8$ span at most 8 dimensions. Section 4 states the condition (p.~315).}
\repair{The Lean statement of Theorem 4.10 assumes that $X$ is connected, and so does the Lean
form of the claim of the introduction.}
\status{Fixed. Blueprint: Theorem 4.10, \lbl{thm:4.10}, and Lemma (a root system containing the
lines), \lbl{lem:4-rootsystem} (\file{blueprint/sec4-leasteigenvalue.tex}). Lean:
\lean{SimpleGraph.hoffman} (\file{LineGraphs/Representation/Hoffman.lean}) and
\lean{SimpleGraph.IsRepresentation.exists_isRepresentedIn_rootType}
(\file{LineGraphs/Representation/RootTypeClassified.lean}), both with the hypothesis
\lean{X.Connected}.}

\errhead{Implicit hypothesis 2}{Theorem 4.13 and ``least eigenvalue $-2$''}{Theorem 4.13; p.~320.
The phrase also occurs in Theorems 4.10 and 4.11 (pp.~319 and 320).}
\printed{``If $G$ is a connected strongly regular graph with least eigenvalue $-2$, then either
\dots''}
\note{Nothing is missing from the printed theorem. Read literally, ``least eigenvalue $-2$'' says
that the least eigenvalue is $-2$, so that $-2$ is an eigenvalue, and the paper uses the phrase in
this sense. Thus Theorem 4.13 is correct as printed. In Theorems 4.10 and 4.11 we read the phrase
as the lower bound $\lambda_{\min}\ge-2$, that is, $2I+A\succeq0$, which makes these two theorems
stronger. Theorem 4.13 is false under that reading.}
\evidence{The paper writes ``the least eigenvalue of $L(G;a_1,\dots,a_m)$ is $-2$ if the number of
vertices exceeds $n$'' (p.~316), and the proof of Theorem 4.13 ends with ``$K_n=L(K_{1,n})$ and
$C_5=L(C_5)$ do not have $-2$ as an eigenvalue''. Under the lower bound reading, $K_n$ with $n>36$
is a counterexample: it is strongly regular in Mathlib's sense (\lean{IsSRGWith}) and has least
eigenvalue $-1$, but it is none of the graphs in (i) or (ii).}
\repair{None is needed. The Lean statement assumes the lower bound and, as a second hypothesis,
that $-2$ is an eigenvalue; together these say that the least eigenvalue is $-2$, by the bridge
lemma \lean{SimpleGraph.leastEigenvalue_eq_neg_two_iff}, so the Lean states the printed
theorem.}
\status{Nothing to fix. Blueprint: Theorem 4.13, \lbl{thm:4.13}, with its two hypotheses
(\file{blueprint/sec4-leasteigenvalue.tex}). Lean: \lean{SimpleGraph.seidel}
(\file{LineGraphs/Representation/Seidel413.lean}), whose hypotheses \lean{hp} and \lean{hneg} say
that $2I+A\succeq0$ and that $-2$ is an eigenvalue; the bridge is
\lean{SimpleGraph.leastEigenvalue_eq_neg_two_iff} (\file{LineGraphs/Spectral/Positivity.lean}).
Theorems 4.10 and 4.11, \lean{SimpleGraph.hoffman} and \lean{SimpleGraph.hoffman_rayChaudhuri},
assume only the lower bound.}

\errhead{Implicit hypothesis 3}{Lemmas 5.4 and 5.5}{Lemmas 5.4 and 5.5; p.~321.}
\printed{Lemma 5.4: ``If $L(G)$ is regular and connected, then either $G$ is regular, or $G$ is
bipartite and the vertices in the same bipartite block have the same valency.'' Lemma 5.5: ``The
spectrum of a regular graph is determined by the spectrum of its line graph.''}
\missing{Lemma 5.4 needs $X$ to have no isolated vertices, and Lemma 5.5 needs at least one edge.}
\evidence{$K_2$ together with an isolated vertex is a counterexample to Lemma 5.4. All edgeless
graphs have the same line graph, so Lemma 5.5 fails without an edge.}
\repair{SP4 of \file{cgss/prerequisites.md}, and the blueprint of Lemmas 5.4 and 5.5, assume
these.}
\status{Fixed. Blueprint: Lemma 5.4, \lbl{lem:5.4}, and Lemma 5.5, \lbl{lem:5.5}
(\file{blueprint/sec5-spectral.tex}). Lean:
\lean{SimpleGraph.regular_or_isBiregularWith_of_lineGraph_regular}
(\file{LineGraphs/Spectral/LineGraphRegular.lean}) assumes that every vertex has positive degree,
and \lean{SimpleGraph.cospectral_of_lineGraph_cospectral}
(\file{LineGraphs/Spectral/LineGraphRecovery.lean}) assumes that both graphs have an edge.}

\errhead{Implicit hypothesis 4}{Hypothesis 2.4 needs two lines}{After Hypothesis 2.4; p.~307.}
\printed{``Without loss of generality we assume that $S$ contains lines $\langle a\rangle$ and
$\langle b\rangle$ whose spanning vectors $a$ and $b$ have the inner product $(a,b)=-1$.''}
\missing{$S$ has at least two lines; in particular, $S$ is not empty.}
\evidence{The assumption is possible exactly when $S$ has at least two lines: an indecomposable set
with two or more lines contains two lines at $60^\circ$, since otherwise all its lines would be
mutually orthogonal, and a choice of signs gives $(\Za,\Zb)=-1$. A single line is an indecomposable
star-closed set, namely $A_1$, and for it $\Gamma_a$, $\Delta$ and $\Sigma$ are not defined. Thus
Lemma 2.12, Theorem 2.13 and Theorem 3.5 need the case of one line separately; the first line of
Theorem 2.13 with $n=1$ is this case. The empty set must be excluded too.}
\repair{The blueprint of Theorems 2.13 and 3.5 splits on the number of lines. In Lean the
structure \lean{LineSystem.Hyp24} carries $\Za$ and $\Zb$, and the classification,
\lean{LineSystem.classification}, assumes a nonempty set (\lean{hne}).}
\status{Fixed. Blueprint: Hypothesis 2.4, \lbl{hyp:2.4}, Lemma 2.12, \lbl{lem:2.12}, step
\lbl{212star}, and Theorem 2.13, \lbl{thm:2.13} (\file{blueprint/sec2-lines.tex}); Theorem 3.5,
\lbl{thm:3.5} (\file{blueprint/sec3-rootsystems.tex}). Lean: \lean{LineSystem.Hyp24}
(\file{LineGraphs/LineSystem/Hyp24.lean}); the split
\lean{LineSystem.IsIndecomposable.eq_pair_or_exists_hyp24}
(\file{LineGraphs/LineSystem/OrthGraphCases.lean}); and \lean{LineSystem.five_possibilities'} and
\lean{LineSystem.classification}, which assume that the set of roots is nonempty.}

\errhead{Implicit hypothesis 5}{Lemma 5.6 needs $X_2$ without isolated vertices}{Lemma 5.6;
p.~321.}
\printed{Lemma 5.6: ``Let $G_1$ be connected bipartite such that $L(G_1)$ is regular. Let $G_2$ be
such that $\operatorname{spec}(L(G_1))=\operatorname{spec}(L(G_2))$. Then either (i) \dots\ or (ii) $G_1$ is
not regular; $G_2$ is regular, not bipartite; $|G_1|=|G_2|+1$.''}
\missing{$X_2$, the paper's $G_2$, has no isolated vertices.}
\evidence{Take $X_1=K_{1,3}$, and let $X_2$ be a triangle together with an isolated vertex. Then
$L(X_1)=L(X_2)=K_3$, but $X_2$ is neither bipartite nor regular. The proof applies Lemma 5.4 to
$X_2$, which needs no isolated vertices (Implicit hypothesis 3); it also uses that $L(X_2)$ is
regular and connected, which follows from Lemmas 5.1 and 5.2. Theorem 5.8 is unaffected, since
there the graph $Y$ with $X_2=L(Y)$ can be chosen without isolated vertices.}
\repair{The blueprint and the Lean statement of Lemma 5.6 assume that $X_2$ has no isolated
vertices.}
\status{Fixed. Blueprint: Lemma 5.6, \lbl{lem:5.6} (\file{blueprint/sec5-spectral.tex}). Lean:
\lean{SimpleGraph.lineGraph_cospectral_bipartite_forall_isBiregularWith}
(\file{LineGraphs/LineGraphSpectra/Bipartite.lean}), whose hypothesis \lean{hnoY} says that every
vertex of the second graph has positive degree.}

\errhead{Implicit hypothesis 6}{The meaning of ``characterized by its spectrum''}{Corollaries
5.9--5.11, Remark 5.12 and Example 5.13; pp.~324--325.}
\printed{For instance, ``The flag graph of a 2-design with $\lambda=1$ is characterized by its
spectrum.''}
\missing{The paper does not define the phrase. Its proofs show, through Theorem 5.8 (ii), that
every graph with the spectrum of $F(D_1)$ is $F(D_2)$ for some 2-design $D_2$ with the same
parameters, and this is the sense of Hoffman and Ray-Chaudhuri [14] and of Doob [6, 7].}
\evidence{Read as ``determined up to isomorphism by its spectrum'', Corollaries 5.10 and 5.11 and
Remark 5.12 are false. The spectrum of $F(D)$ depends only on the parameters, since
$NN^T=(r-\lambda)I+\lambda J$ (proof of Theorem 5.8), whereas $F(D)=L(I(D))$ determines $I(D)$ by
Whitney's theorem, and so determines $D$ up to duality. Thus the two Steiner triple systems on 13
points have cospectral, non-isomorphic flag graphs (against Corollary 5.11), and the five symmetric
designs with parameters $(15,7,3)$ give at least three (against Corollary 5.10 and Remark 5.12,
with $\lambda=3$). For Corollary 5.9 the two readings agree, since the trivial design with given
parameters is unique.}
\repair{The blueprint and the Lean statements define the phrase in this sense: every graph with the
spectrum of $F(D_1)$ is isomorphic to $F(D_2)$ for some 2-design $D_2$ with the parameters of
$D_1$. On the paper's route, the ``not characterized'' of Example 5.13 is then EX4 of
\file{cgss/prerequisites.md}, a case of Whitney's theorem; the Lean proof avoids EX4 by a clique
bound (blueprint, Example 5.13).}
\status{Fixed. Blueprint: Definition (incidence graphs, flag graphs, characterization),
\lbl{def:flag}, used from Corollary 5.9, \lbl{cor:5.9}, to Example 5.13, \lbl{ex:5.13}
(\file{blueprint/sec5-spectral.tex}). Lean: \lean{Design.FlagGraphCharacterized}
(\file{LineGraphs/LineGraphSpectra/DesignGraphs.lean}), used by
\lean{Design.IsSymmetric.flagGraphCharacterized}, \lean{Design.flagGraphCharacterized_of_lam_eq_one},
\lean{Design.flagGraphCharacterized_of_lam_lt_nineteen} and
\lean{Design.not_flagGraphCharacterized_of_resolvable}, the last for Example 5.13.}

\errhead{Implicit hypothesis 7}{Definition 2.2 means maximal within the span}{Section 2, first
paragraph, and Definition 2.2; p.~306.}
\printed{``Without loss of generality the lines are taken to pass through the origin of $\re^n$
and to span $\re^n$.'' and ``A set of lines at $60^\circ$ and $90^\circ$ in $\re^n$ is maximal if
any line of $\re^n$ belongs to the set whenever it has angles $60^\circ$ or $90^\circ$ with all
lines of the set.''}
\missing{Maximality is meant relative to the space that the lines span.}
\evidence{Theorem 3.6 needs this reading: it calls $A_7$ and $A_8$ non-maximal because
$A_7\subset E_7$ and $A_8\subset E_8$ in the same dimension, and $A_n$ is defined in
Definition 3.1 in $\re^{n+1}$, where $\Ze_1+\Ze_2$ is at $60^\circ$ or $90^\circ$ to all its lines,
so that no $A_n$ would be maximal in its ambient space. Brouwer, Cohen and Neumaier,
\emph{Distance-Regular Graphs}, work with root lattices and do not use the notion.}
\repair{The Lean definition and Definition 2.2 of the blueprint ask for maximality within
$\operatorname{span}(\Phi)$. Lemma 2.3 is unchanged, since the sum of two roots lies in
the span.}
\status{Fixed. Blueprint: Definition 2.2, \lbl{def:2.2}, and Lemma 2.3, \lbl{lem:2.3}
(\file{blueprint/sec2-lines.tex}); Theorem 3.6, \lbl{thm:3.6}
(\file{blueprint/sec3-rootsystems.tex}). Lean: \lean{LineSystem.IsMaximal}
(\file{LineGraphs/LineSystem/Defs.lean}), with \lean{LineSystem.IsMaximal.isStarClosed} for
Lemma 2.3 and \lean{LineSystem.classification_of_isMaximal} for Theorem 3.6.}

\errhead{Implicit hypothesis 8}{Theorem 3.12: a part may be empty}{Theorem 3.12; p.~315.}
\printed{The lines can be ``partitioned into two subsets''.}
\missing{Either subset may be empty.}
\evidence{With both subsets nonempty the theorem is false: a single line lies in a Root System but
cannot be split into two nonempty parts, and two lines of $A_2$ at $60^\circ$ must lie in the same
part.}
\repair{The partition is a colouring $\phi$ of $S$ with two colours, not necessarily surjective,
and the empty set of lines is allowed. The note \file{cgss/proofs/thm312/thm312.tex} proves the theorem
in this form.}
\status{Fixed. Blueprint: Theorem 3.12, \lbl{thm:3.12}, with its paragraph on the conventions
(\file{blueprint/sec3-thm312.tex}), after \file{cgss/proofs/thm312/thm312.tex}. Lean:
\lean{LineSystem.subset_rootSet_iff_bipartition} (\file{LineGraphs/Theorem312/Main.lean}), whose
colouring, in \lean{LineSystem.HasBipartition}, takes values in \lean{Bool}; its docstring says
that the set and either colour fibre may be empty.}

\end{document}
